% ==================================================================
\section{模型评价}\label{sec:eval}
% ==================================================================

\subsection{模型的优点}

\begin{enumerate}[label=\arabic*.]
\item \textbf{特征位置都能反查。}问题一用强制对齐给定转写得到词级时间，每个有效位置都挂着“原文字符—词—子词位置—音视频帧”的反查链；典型样本逐词对得上秒数和帧数，31段583个词的人工核验全部通过。
\item \textbf{失败被如实保留。}100条样本一条不删，拿不到时间的位置不补造时间，提取失败的模态记为 None 而不是零，8段没有词时间的样本、4段失败样本都在台账中可查。
\item \textbf{三种零值分开处理。}问题二把补齐填充、原生不可观测、人为遮蔽分成三个掩码，文本遮蔽走重新编码而不是置零，模型和评估都不会把“没有”误读成“为零”。
\item \textbf{简单模型胜在稳定。}PM-B2 在7个基线和4个时序网络对照中缺失面板 F1 最高，三个种子间的标准差只有0.0013；比较时按源视频配对重抽样，没有把重复出现的样本当成独立样本。
\item \textbf{解释分层检验。}问题三的 Shapley 值精确计算且逐样本满足效率恒等式，局部证据有高分、低分、随机三组对照，时间证据再经人工核验，三层检验各管一件事。
\end{enumerate}

\subsection{模型的缺点}

\begin{enumerate}[label=\arabic*.]
\item \textbf{问题一的特征与附件2不通用。}自提的声学特征是25维 eGeMAPS，附件2是74维；两套特征不能互换，问题一的特征质量也就没法用问题二的模型直接检验。
\item \textbf{中性类识别较差。}测试集中性召回只有0.361、F1 只有0.397；B 阶段的加权损失和层次头提高了召回，却降低了精确率，没有同时改善。
\item \textbf{改进实验只有一个种子。}A–E 的15种配置都只跑了 seed 17，胜出者没有复核，蒸馏中关系保持项的单独作用也没有做去除消融。
\item \textbf{专项测试集没有标签。}附件3、4的可信度只能借助附件2上的表现间接判断，附件4的删除检验只有20条样本。
\item \textbf{缺失面板是预设的干预。}63档局部面板和压力面板都是我们设定的缺失方式，真实场景中的缺失未必如此，结论不能直接推广到任意缺失机制。
\end{enumerate}

\subsection{改进方向}

第一，对 A–E 中真正可能改变结论的候选（如 LH-B2、D2）补跑 seed 29 和 43，并对 E2 的关系保持项做单独去除，再用按源视频划分的训练集内交叉验证检验选型是否稳定；在这之前，不用已经看过的测试集或无标签专项集反过来调参。第二，针对中性类，可以把分类阈值和强度预测联合校准，例如让强度接近0且分类置信度低的样本优先判为中性，并在验证集上确定阈值。第三，问题一可以把自提特征接到一个在附件2上训练、但输入维度与之匹配的模型上，才能真正比较不同对齐方式对下游预测的影响。第四，对非文本主导且置信度低的样本，可以设计自动转交人工复核的规则，把问题三的解释直接用到实际流程里。

% ==================================================================
\section*{人工智能工具使用说明}
\addcontentsline{toc}{section}{人工智能工具使用说明}
% ==================================================================

本文的写作过程使用了大语言模型工具：原稿的论文框架、文字初稿及部分图表生成脚本借助 \todo{按实际会话记录填写工具名称与型号，原稿写的是 OpenAI GPT 系列（Codex 会话）}完成；本稿的结构调整、文字改写与 LaTeX 排版借助 Anthropic Claude 完成。文中所有指标、表格和图均来自队伍保存的实验结果文件、冻结清单和人工核验记录，改写过程中没有新增或改动任何实验数值；新增的两张图（图\ref{fig:law}、图\ref{fig:special}）只是把原有表格中的数字重新画成图。参赛队已对照原始记录逐项复核上述内容\todo{复核完成后删去此提示}。若竞赛系统要求提供提示词、模型输出及后续处理方式，另按要求提交。
